% =====================================================================
\subsection{Fase 4: inteligencia artificial}

\subsubsection{Modelo de los androides SENTINEL}

Los androides de seguridad se construyeron sobre el modelo \textit{Y Bot} de Mixamo
\parencite{mixamo2026}, configurado como Avatar Humanoid; por eso usan las mismas animaciones de
reposo, caminar y correr que los técnicos. Una herramienta de editor les da su aspecto: cuerpo de
metal negro pulido, juntas que brillan (en la oscuridad se ve la silueta antes que el cuerpo), un
visor rojo pegado al hueso de la cabeza y una linterna roja que marca hacia dónde miran
(Figura~\ref{fig:androide}).

\FiguraDoble{042_androide-sentinel-primer-plano.jpg}{043_androide-detecta-y-persigue.jpg}{fig:androide}{Androide SENTINEL (izquierda) y un androide que detecta y persigue al técnico (derecha)}{}

\subsubsection{Percepción}

\begin{itemize}
  \item \textbf{Visión:} un cono invisible de 9~m de alcance y 55° de apertura, hecho con un
        \textit{Mesh Collider} convexo en modo \textit{trigger}, avisa cuando el técnico entra en él.
        Entonces el androide lanza un rayo desde sus ojos hasta el pecho del técnico: si choca antes
        con un rack, una pared o un mueble alto, no lo ve. Por eso esconderse detrás de los racks
        funciona. A menos de 1,6~m lo percibe aunque esté a su espalda.
  \item \textbf{Oído:} el androide va a investigar si el técnico corre a menos de 5~m, si dispara
        cerca (18~m la pistola, 24~m la escopeta), si una máquina expendedora suelta un producto o si
        una válvula se abre fuera de orden. Caminar es la forma sigilosa de moverse.
\end{itemize}

\subsubsection{Máquina de estados}

El comportamiento de cada androide está en el script \cod{AndroideIA}, organizado como una máquina
de estados (Tabla~\ref{tab:estados}). Los androides se mueven con un \textit{NavMeshAgent} sobre
una malla de navegación que se calcula al cargar cada nivel a partir del suelo, las paredes, las
filas de racks y los muebles \parencite{unitymanual2026}; las puertas cerradas llevan un obstáculo
que corta la malla hasta que se abren. Patrullan a 1,3~m/s y persiguen a unos 3,5~m/s, más despacio
que el técnico cuando corre (5~m/s), de modo que huir siempre es posible.

\begin{TablaAPA}{tab:estados}{Estados de la inteligencia artificial de los androides}
\begin{tabularx}{\linewidth}{P{2.6cm}LP{4.4cm}}
\toprule
Estado & Comportamiento & Pasa a\ldots \\
\midrule
Reposo & Anclado en su estación de carga, con el visor casi apagado y un pulso tenue & Búsqueda al sonar la alarma o si el técnico pasa muy cerca \\
Patrulla & Recorre sus puntos de ruta a paso lento y se detiene un momento en cada uno & Persecución si ve al técnico; Búsqueda si lo oye \\
Persecución & Corre hacia el técnico; la linterna se intensifica y parpadea & Recalibrando si lo alcanza; Búsqueda si lo pierde de vista 3~s \\
Búsqueda & Va al último lugar donde vio u oyó al técnico y mira alrededor & Persecución si lo ve; Patrulla si no lo encuentra \\
Recalibrando & Se detiene 2,5~s después de alcanzar al técnico, dándole tiempo para huir & Patrulla \\
Atraído & Va al señuelo o al ruido y se queda ahí mientras dure & Búsqueda o Persecución \\
Apagado & Desplomado por el KILLSWITCH; recibe el doble de daño & Patrulla, después de 10~s \\
\bottomrule
\end{tabularx}
\end{TablaAPA}

\subsubsection{Estaciones de carga y alarma}

Algunos androides esperan en estaciones de carga: un panel, una base con luz y la unidad anclada en
reposo (Figura~\ref{fig:reposo}). Se reactivan si el técnico pasa pegado a ellos, corre cerca o se
les planta enfrente, y todos despiertan cuando se extrae el Disco Maestro de un nivel. En ese
momento suena la alarma, las luces pasan a rojo y, durante 22~s, el sistema les comunica la
posición del técnico aunque no lo vean, lo que obliga a huir.

\FiguraDoble{073_androide-en-reposo.jpg}{097_alarma-androide-se-enciende.jpg}{fig:reposo}{Androide en reposo (izquierda) y una unidad que se enciende en su estación de carga al sonar la alarma (derecha)}{}

\subsubsection{Jefe final: SENTINEL-\texorpdfstring{$\Omega$}{Omega}}

En el Nivel 3, al tomar el último disco, CERBERUS se descarga en una unidad prototipo que esperaba en
una cápsula junto a la salida de emergencia. Es un androide 1,45 veces más grande, con juntas
doradas y 60 puntos de vida, que bloquea la salida hasta ser destruido. Su script
\cod{JefeSentinel} sigue un ciclo de tres fases (Tabla~\ref{tab:jefe}).

\begin{TablaAPA}{tab:jefe}{Ciclo de combate del jefe final}
\begin{tabularx}{\linewidth}{P{3.6cm}P{1.6cm}P{2.6cm}L}
\toprule
Fase & Duración & Daño que recibe & Qué debe hacer el jugador \\
\midrule
Blindado: persigue al técnico & 8 s & La mitad & Alejarse y no gastar balas \\
Cargando una descarga & 2,2 s & Normal & Alejarse: la descarga hiere a menos de 6,5~m \\
Expuesto: quieto y sobrecalentado & 2,5 s & El doble & Atacarlo, incluso con la bayoneta \\
\bottomrule
\end{tabularx}
\end{TablaAPA}

A la mitad de su vida llama a dos unidades de refuerzo que estaban en reserva y carga más seguido.
Arriba de la pantalla se muestra una barra de vida en texto con su estado (blindado, cargando o
expuesto). Al destruirlo, CERBERUS pierde la unidad que lo aloja: todos los androides se desconectan
y la salida se desbloquea (Figuras~\ref{fig:jefe1} y~\ref{fig:jefe2}). El combate puede ganarse
incluso sin munición, aprovechando con la bayoneta las ventanas en que queda expuesto.

\FiguraDoble{101_jefe-sentinel-omega.jpg}{102_jefe-cargando-descarga.jpg}{fig:jefe1}{SENTINEL-$\Omega$ al activarse (izquierda) y cargando su descarga (derecha)}{}

\FiguraDoble{103_jefe-expuesto.jpg}{105_jefe-destruido.jpg}{fig:jefe2}{SENTINEL-$\Omega$ expuesto (izquierda) y destruido (derecha)}{}

% =====================================================================
\subsection{Fase 5: audio}

\subsubsection{Selección de la música y los efectos}

La música y los efectos se tomaron de Pixabay, cuya licencia permite usarlos de forma gratuita,
también con fines comerciales y sin obligación de atribución \parencite{pixabaylicencia}; aun así,
cada autor se acredita en el juego y en el Anexo~B. Se buscó una atmósfera de tensión constante,
como la de los survival horror de referencia. La Tabla~\ref{tab:musica} resume las pistas, y los
30 efectos cubren las armas, los androides, el técnico, la interfaz, el entorno, el jefe y el
ascensor.

\begin{TablaAPA}{tab:musica}{Música y ambientes del juego}
\begin{tabularx}{\linewidth}{P{4.4cm}L}
\toprule
Momento del juego & Pista \\
\midrule
Menú principal & \textit{Horror piano} \parencite{leberch_piano} \\
Exploración de los niveles & \textit{Survival horror} \parencite{leberch_survival} \\
Persecución & \textit{Suspense horror tension} \parencite{delosound_tension} \\
Combate contra el jefe & \textit{Dark fight music boss III} \parencite{montogoronto_boss} \\
Pantalla de Victoria & \textit{Victory} \parencite{paulyudin_victory} \\
Ambiente del centro de datos & \textit{Server room atmos} \parencite{freesound_servidores} \\
Ascensor entre niveles & \textit{080223-002 freight elevator} \parencite{freesound_ascensor} \\
\bottomrule
\end{tabularx}
\end{TablaAPA}

\subsubsection{Preparación con FFmpeg}

Los archivos se prepararon con FFmpeg \parencite{ffmpeg2026}, mediante un script que se puede
volver a ejecutar en cualquier momento:

\begin{enumerate}
  \item Se midió en cada grabación dónde hay sonido y dónde silencio, y su volumen cada 100~ms,
        para elegir el tramo útil. Por ejemplo, de una grabación de 11~s con varios ciclos de
        escopeta se tomó un solo bombeo, y de una grabación de pasos se sacaron cuatro pisadas
        sueltas que el juego alterna al azar.
  \item Cada efecto se recortó con fundidos de pocos milisegundos para que no produzca clics.
  \item Cada sonido se niveló a una sonoridad (LUFS) acorde con su papel en la mezcla: pasos
        suaves, disparos fuertes, interfaz discreta.
  \item La música y los ambientes se convirtieron en bucles sin corte: el final de la pista se
        funde con su comienzo, de modo que al repetirse no se nota el salto.
  \item Los efectos se guardaron en WAV mono y la música en OGG estéreo.
\end{enumerate}

\Figura[0.9\linewidth]{112_audio-pixabay-recortes.jpg}{fig:recortes}{Grabaciones originales y su versión recortada y nivelada para el juego}{}

\subsubsection{Banco de sonidos y gestor de audio}

Los sonidos se reúnen en un \textit{banco de sonidos} (un archivo de datos de Unity) que una
herramienta de editor llena y configura: la música y el ambiente se leen del disco mientras suenan
(\textit{streaming}), y los efectos se cargan descomprimidos y en mono para sonar al instante y en
3D. El banco admite varias tomas de un mismo efecto, como las pisadas, y un quejido distinto para
cada personaje al recibir daño.

El gestor \cod{AudioNexus} se crea solo al iniciar el juego y pasa de una escena a otra. Mezcla la
música por capas, cuyos volúmenes se funden según lo que ocurre:

\begin{itemize}
  \item en el menú suena su tema;
  \item en los niveles suenan la música de exploración y el zumbido de la sala de servidores;
  \item la capa de \textbf{tensión} entra cuando algún androide persigue al técnico;
  \item la \textbf{sirena} suena mientras dura la alarma;
  \item en el combate contra el jefe entra su música y baja la de exploración;
  \item dentro del ascensor la mezcla se apaga y la alarma se oye lejos;
  \item al final suena la pieza de Victoria o el remate del Game Over.
\end{itemize}

Los efectos suenan en 3D desde donde ocurren, con un grupo de 20 fuentes reutilizables y una leve
variación de tono para que los sonidos repetidos no suenen idénticos. Una separación mínima entre
sonidos del mismo tipo evita que, por ejemplo, seis androides que detectan al técnico a la vez
produzcan seis golpes encimados. El \textit{oído} del juego (\textit{AudioListener}) va en el
técnico y no en la cámara: con cámaras fijas que cortan de golpe, el sonido saltaría de un lado a
otro con cada corte, mientras que con el oído en el personaje un androide que se acerca por la
izquierda siempre suena a su izquierda.
